\uuid{3tcx}
\exo7id{1899}
\titre{exo7 1899}
\auteur{legall}
\organisation{exo7}
\datecreate{2003-10-01}
\isIndication{false}
\isCorrection{false}
\chapitre{Espace vectoriel normé}
\sousChapitre{Espace vectoriel normé}
\cpge{mpsi-mp2i.calcul-matriciel-systemes.e, pcsi.determinants.c, ptsi.determinants.a, mpsi-mp2i.determinants.d, pcsi.calcul-matriciel-systemes-lineaires.d, ptsi.calcul-matriciel-systemes-lineaires.d, bcpst1.matrices, mp-mpi.reduction-endomorphismes.f, mp-mpi.topologie-espaces-normes.b, pc.espaces-vectoriels-normes.b, psi.espaces-vectoriels-normes.b, bcpst2.valeurs-propres-vecteurs-propres.a}
\module{Topologie}
\niveau{L3}
\difficulte{}

\contenu{
\begin{enumerate}
    \item \question{Soit $A\in M_n(\Cc )$. Montrer quil existe une suite de 
matrices $(A_n)_{n\in \Nn }$
inversibles convergeant vers $A$ (en un sens que l'on pr\'ecisera).}
    \item \question{Soit $N\in M_n(\Cc )$ une matrice nilpotente. Calculer les 
valeurs propres de $N$. Montrer
que $\hbox{det}(I+N)=1.$}
    \item \question{Soit $A\in M_n(\Cc )$ telle que $AN=NA$. Calculer $\hbox{det}(A+N).$}
\end{enumerate}
}
